\documentclass[12pt]{article}
\usepackage[utf8]{inputenc}
\usepackage[portuguese]{babel}
\usepackage{amsmath}
\usepackage{siunitx}
\usepackage{graphicx}
\usepackage{geometry}
\usepackage{multicol}
\usepackage{circuitikz}
\usepackage{parskip}
\usepackage{enumitem}

\geometry{a4paper, left=1cm, right=1.5cm, top=1cm, bottom=1.5cm}

\newcommand{\cabecalho}{
	\hfill
	\begin{minipage}[t]{0.7\textwidth}
		\raggedleft
		\textbf{Teste Modelo 1, Fundamentos de Física}
	\end{minipage}
	\vspace{0.5cm}
}

\newcommand{\pontos}[1]{\hfill\textbf{(#1 valores)}}

\title{Teste Modelo 1}
\author{}
\date{20 de março de 2026}

\setlength{\parindent}{0pt}
\setlist[enumerate]{itemsep=4pt, topsep=4pt}

\begin{document}
	
	\cabecalho
	
	\textbf{Exercício 1}. Considere uma piscina cilíndrica com base circular de raio $2\,\text{m}$ e altura $3\,\text{m}$. A piscina está cheia de um líquido até $30\,\text{cm}$ da borda.
	\begin{enumerate}[label=\textbf{\alph*)}]
		\item Calcule o volume total da piscina. \pontos{1,5}
		\item Sabendo que o líquido tem uma densidade $\rho = 600\,\text{kg/m}^3$, qual é a massa do líquido na piscina? \pontos{1,0}
	\end{enumerate}
	
	\vspace{0.4cm}
	
	\textbf{Exercício 2}. Considere o vetor $\overline{v} = (v_x, v_y)$ de módulo 2 e orientação definida pelo ângulo $\theta = 60^\circ$.
	\begin{enumerate}[label=\textbf{\alph*)}]
		\item Desenhe o vetor no plano cartesiano. \pontos{0,5}
		\item Calcule as suas componentes e escreva $\overline{v}$ com os vetores unitários $\overline{i}$ e $\overline{j}$. \pontos{1,5}
		\item Calcule o vetor unitário com a mesma orientação de $-\overline{v}$. \pontos{1,0}
		\item Se $\overline{w} = (0; 1,3)$, quais são os valores de $|\overline{w}|$ e $\phi$, tais que $w_x = |\overline{w}| \cos \phi$ e $w_y = |\overline{w}| \sin \phi$? \pontos{1,5}
	\end{enumerate}
	
	\vspace{0.4cm}
	
	\textbf{Exercício 3}. Uma partícula no plano descreve um movimento acelerado (não necessariamente uniformemente acelerado), passando de um ponto na posição $\overline{r}_A = (1; 1)\,\text{m}$ para a posição $\overline{r}_B = (1; 1,5)\,\text{m}$, percorrendo uma distância $d = 65\,\text{cm}$.
	\begin{enumerate}[label=\textbf{\alph*)}]
		\item O movimento é retilíneo ou curvilíneo? Desenhe uma possível trajetória que passe por A e B. \pontos{1,0}
		\item Devido a um mecanismo, no ponto B a velocidade da partícula passa de $\overline{v}_1 = (2, -1)\,\text{m/s}$ para uma velocidade $\overline{v}_2 = (0, -1)\,\text{m/s}$ num intervalo de tempo de $\Delta t = 750\,\text{ms}$. Calcule a aceleração média da partícula durante $\Delta t$ e explique o significado do seu sinal. \pontos{2,5}
		\item Depois do intervalo $\Delta t$, o corpo continua em movimento retilíneo uniformemente variado, sem mudar a orientação da trajetória, e imobiliza-se após um intervalo $\Delta T = 15\,\text{s}$. Qual é a aceleração durante esse intervalo? \pontos{1,5}
		\item Desenhe a trajetória durante $\Delta T$ e calcule a distância percorrida. \pontos{1,5}
	\end{enumerate}
	
	\vspace{0.4cm}
	
	\textbf{Exercício 4}. Uma bola é lançada do topo de um prédio de altura $30 \text{ m}$. No instante do lançamento, a bola tem uma velocidade de $10 \text{ m/s}$, apontando para baixo e formando um ângulo de $45,0^\circ$ com a vertical. Desprezando a resistência do ar, determine:
	
	\vspace{0.15cm}
	\begin{enumerate}[label=\textbf{\alph*)}]
		\item Desenhe a trajetória do projétil, explicitando numericamente o ângulo de lançamento. \pontos{0,5}
		\item Qual é a altura máxima atingida durante a trajetória? \pontos{1,0}
		\item Qual é o instante em que o projétil atinge o solo? \pontos{2,0}
		\item A que distância horizontal do prédio o projétil atinge o solo? \pontos{1,0}
		\item Calcule o ângulo de impacto e a velocidade no ponto de impacto, desenhando o seu sentido e orientação. \pontos{2,0}
	\end{enumerate}
	
\end{document}